\documentclass[12pt, titlepage]{article}

\usepackage{amsmath, mathtools}

\usepackage[round]{natbib}
\usepackage{amsfonts}
\usepackage{amssymb}
\usepackage{graphicx}
\usepackage{colortbl}
\usepackage{xr}
\usepackage{hyperref}
\usepackage{longtable}
\usepackage{xfrac}
\usepackage{tabularx}
\usepackage{float}
\usepackage{siunitx}
\usepackage{booktabs}
\usepackage{multirow}
\usepackage[section]{placeins}
\usepackage{caption}
\usepackage{fullpage}

\hypersetup{
bookmarks=true,     % show bookmarks bar?
colorlinks=true,       % false: boxed links; true: colored links
linkcolor=red,          % color of internal links (change box color with linkbordercolor)
citecolor=blue,      % color of links to bibliography
filecolor=magenta,  % color of file links
urlcolor=cyan          % color of external links
}

\usepackage{array}

\externaldocument{../../SRS/SRS}

\input{../../Comments.text}
\input{../../Common.text}

\begin{document}

\title{Module Interface Specification for \progname{}}

\author{\authname}

\date{\today}

\maketitle

\pagenumbering{roman}

\section{Revision History}

\begin{tabularx}{\textwidth}{p{3cm}p{2cm}X}
\toprule {\bf Date} & {\bf Version} & {\bf Notes}\\
\midrule
2026/03/25 & 1.0 & Initial MIS\\
2026/04/07 & 1.1 & Modified based on the reviewer's feedback\\
\bottomrule
\end{tabularx}

~\newpage

\section{Symbols, Abbreviations and Acronyms}

See SRS Documentation at \href{https://mei0725.github.io/FlockingSimulator/SRS/SRS.pdf}{SRS.pdf}

\newpage

\tableofcontents

\newpage

\pagenumbering{arabic}

\section{Introduction}

The following document details the Module Interface Specifications for
\progname \, which simulates flocking behaviors by different ODE methods, 
visualizes animation of result, and analyzes parameter effects through charts.

Complementary documents include the System Requirement Specifications
and Module Guide. The full documentation and implementation can be
found at \href{https://mei0725.github.io/FlockingSimulator/}{here}

\section{Notation}

The structure of the MIS for modules comes from \citet{HoffmanAndStrooper1995},
with the addition that template modules have been adapted from
\cite{GhezziEtAl2003}.  The mathematical notation comes from Chapter 3 of
\citet{HoffmanAndStrooper1995}.  For instance, the symbol := is used for a
multiple assignment statement and conditional rules follow the form $(c_1
\Rightarrow r_1 | c_2 \Rightarrow r_2 | ... | c_n \Rightarrow r_n )$.

The following table summarizes the primitive data types used by \progname. 

\begin{center}
\renewcommand{\arraystretch}{1.2}
\noindent 
\begin{tabular}{l l p{7.5cm}} 
\toprule 
\textbf{Data Type} & \textbf{Notation} & \textbf{Description}\\ 
\midrule
character & char & a single symbol or digit\\
integer & $\mathbb{Z}$ & a number without a fractional component in (-$\infty$, $\infty$) \\
natural number & $\mathbb{N}$ & a number without a fractional component in [1, $\infty$) \\
real & $\mathbb{R}$ & any number in (-$\infty$, $\infty$)\\
\bottomrule
\end{tabular} 
\end{center}

\noindent
The specification of \progname \ uses some derived data types: sequences, strings, and
tuples. Sequences are lists filled with elements of the same data type. Strings
are sequences of characters. Tuples contain a list of values, potentially of
different types. In addition, \progname \ uses functions, which
are defined by the data types of their inputs and outputs. Local functions are
described by giving their type signature followed by their specification.

\section{Module Decomposition}

The following table is taken directly from the Module Guide document for this project.

\begin{table}[h!]
\centering
\begin{tabular}{p{0.3\textwidth} p{0.6\textwidth}}
\toprule
\textbf{Level 1} & \textbf{Level 2}\\
\midrule

{Hardware-Hiding} & ~ \\
\midrule

\multirow{7}{0.3\textwidth}{Behaviour-Hiding Module}
& Main Controller Module \\
& Input Validation Module\\
& Force Module \\
& ODE Solver Module\\
& Flocking Metrics Module\\
& Dynamic Module \\
& Flocking Parameter Module \\
& Animation Status Module \\
& Flocking Status Module \\
\midrule

\multirow{3}{0.3\textwidth}{Software Decision Module}
& Input Reader Module\\
& File I/O Module\\
& Animation Simulator Module\\ 
& Chart Drawer Module\\
\bottomrule
\end{tabular}
\caption{Module Hierarchy}
\label{TblMH} 
\end{table}

\newpage
~\newpage

\section{MIS of Main Controller Module} \label{mMC}

\subsection{Module}

MainController

\subsection{Uses}

FilesIO, InputReader, Dynamic, AnimationSimulator, ChartDrawer

\subsection{Syntax}

\subsubsection{Exported Constants}

\begin{itemize}
\item $startRadius$ := 5.
\end{itemize}

\subsubsection{Exported Access Programs}

\begin{center}
\begin{tabular}{p{4cm} p{4cm} p{4cm} p{2cm}}
\hline
\textbf{Name} & \textbf{In} & \textbf{Out} & \textbf{Exceptions} \\
\hline
\texttt{InitClicked} & - & - & - \\
\texttt{CalculateClicked} & - & - & - \\
\texttt{PlayAnimaClicked} & - & - & - \\
\texttt{DrawClicked} & - & - & - \\
\hline
\end{tabular}
\end{center}

\subsection{Semantics}

\subsubsection{State Variables}

None.

\subsubsection{Access Routine Semantics}

\noindent InitClicked():
\begin{itemize}
\item transition: Read the input parameters from Input Reader Module, add the 
initial positions of the flocking simulation, and save it using File I/O Module 
by following steps. 
\end{itemize}
\textit{\# Get user input from InputReader and set it in the FlockingParameter object}

parameters :=  GetUserInput()

\textit{\# Initialize agents' starting positions}

FlockingInit(parameters)

\textit{\# Save the parameters to a file by FilesIO}

path := Save(parameters)

\textit{\# Show the path in GUI}

\vspace{2em}

\noindent CalculateClicked():

\begin{itemize}
\item transition: Read the input parameters from the path users input from 
Input Reader Module, calculate the process of simulation using Dynamics Module, 
and save the results using File I/O Module by following steps. 
\end{itemize}
\textit{\# Get the FlockingParameter parameters and method from InputReader}

parameters :=  GetSimulateParam()

method := GetMethod()

\textit{\# Calculate the metrics of flocking in the simulation by Dynamics}

[animationStatus, flockMetrics] := FlockSimulation(parameters, method)

\textit{\# Save the parameters to a file by FilesIO}

animaPath := Save(animationStatus)

flockPath := Save(flockMetrics)

\textit{\# Show the path in GUI}

\vspace{2em}

\noindent PlayAnimaClicked():

\begin{itemize}
\item transition: Read the animation information from the path users input from 
Input Reader Module, generate the animation and display it using Animation 
Simulator Module by following steps. 
\end{itemize}
\textit{\# Get the AnimationStatus anime from InputReader}

anima :=  GetAnimaStatus()

\textit{\# Generate the animation and display by AnimationSimulator}

PlaySimulation(anima)

\vspace{2em}

\noindent DrawClicked():

\begin{itemize}
\item transition: Read the animation information from the path users input from 
Input Reader Module, draw the line charts and display it by following steps. 
\end{itemize}
\textit{\# Get the FlockingStatus flock from InputReader}

flock :=  GetFlockStatus()

\textit{\# Draw the line charts and display by ChartDrawer}

DrawChart(flock)

\subsubsection{Local Functions}
\noindent FlockingInit():
\begin{itemize}
\item input: \texttt{parameters}: FlockingParameter
\item output: None
\item description: Generate the initial positions for a specified number of 
agents and assign them to \texttt{parameter.starts}.
\end{itemize}

~\newpage

\section{MIS of Input Validation Module} \label{mIV}

\subsection{Module}

InputValidation

\subsection{Uses}

N/A

\subsection{Syntax}

\subsubsection{Exported Constants}

\begin{itemize}
\item $k_{min}$ := 0.
\item $k_{max}$ := 2.0.
\item $N_{min}$ := 1.
\item $N_{max}$ := 5000.
\item $t_{min}$ := 0.005.
\item $t_{max}$ := 0.1.
\item $p_{min}$ := -500.
\item $p_{max}$ := 1000.
\item $r_{min}$ := 1.
\item $r_{max}$ := 500.
\item $EEULER$ := EEuler.
\item $RK4$ := RK4.
\item $ISEULER$ := ISEuler.
\end{itemize}

\subsubsection{Exported Access Programs}

\begin{center}
\begin{tabular}{p{4cm} p{4cm} p{4cm} p{2cm}}
\hline
\textbf{Name} & \textbf{In} & \textbf{Out} & \textbf{Exceptions} \\
\hline
\texttt{KVerify} & \texttt{input}: str, \texttt{type}: type & - & MissingK, ErrorK \\
\texttt{TVerify} & \texttt{input}: str & - & MissingT, ErrorT \\
\texttt{NVerify} & \texttt{input}: str & - & MissingN, ErrorN \\
\texttt{ObsVerify} & \texttt{input}: str & - & MissingP, MissingR, ErrorP, ErrorR \\
\texttt{MethodVerify} & \texttt{method}: str & - & MissingMethod, ErrorMethod \\
\texttt{PathVerify} & \texttt{path}: str & - & MissingPath \\
\hline
\end{tabular}
\end{center}

\subsection{Semantics}

\subsubsection{Access Routine Semantics}

\noindent KVerify(input, type):
\begin{itemize}
\item output: out := None. 
\item exception: exc :=
\begin{itemize}
  \item MissingK, if $input$ is null.
  \item ErrorK, if $input$ is not number or $input < k_{min}$ or $input > k_{max}$.
\end{itemize}
\end{itemize}

\noindent TVerify(input):
\begin{itemize}
\item output: out := None. 
\item exception: exc :=
\begin{itemize}
  \item MissingT, if $input$ is null.
  \item ErrorT, if $input$ is not number or $input < t_{min}$ or $input > t_{max}$.
\end{itemize}
\end{itemize}

\noindent NVerify(input):
\begin{itemize}
\item output: out := None. 
\item exception: exc :=
\begin{itemize}
  \item MissingN, if $input$ is null.
  \item ErrorN, if $input$ is not number or $input < N_{min}$ or $input > N_{max}$.
\end{itemize}
\end{itemize}

\noindent ObsVerify(input):
\begin{itemize}
\item output: out := None. 
\item exception: exc :=
\begin{itemize}
  \item MissingP, if for any item $i$ in $input$, $input[i].p$ is null.
  \item MissingR, if for any item $i$ in $input$, $input[i].r$ is null.
  \item ErrorP, if for any item $i$ in $input$, $input[i].p.x < p_{min}$ 
  or $input[i].p.x > p_{max}$ or $input[i].p.y < p_{min}$ 
  or $input[i].p.y > p_{max}$.
  \item ErrorR, if for any item $i$ in $input$, $input[i].r < r_{min}$ 
  or $input[i].r > r_{max}$.
\end{itemize}
\end{itemize}

\noindent MethodVerify(method):
\begin{itemize}
\item output: out := None. 
\item exception: exc :=
\begin{itemize}
  \item MissingMethod, if $method$ is null.
  \item ErrorMethod, if $method$ is not included in {$EEULER$, $RK4$, $ISEULER$}
\end{itemize}
\end{itemize}

\noindent PathVerify(path):
\begin{itemize}
\item output: out := None. 
\item exception: exc :=
\begin{itemize}
  \item MissingPath, if $path$ is null.
\end{itemize}
\end{itemize}

~\newpage

\section{MIS of Force Module} \label{mForce}

\subsection{Module}

Force

\subsection{Uses}

N/A

\subsection{Syntax}

\subsubsection{Exported Constants}

\begin{itemize}
\item $v_{max}$ := 14.
\end{itemize}

\subsubsection{Exported Access Programs}

\begin{center}
\begin{tabular}{p{4cm} p{4cm} p{4cm} p{2cm}}
\hline
\textbf{Name} & \textbf{In} & \textbf{Out} & \textbf{Exceptions} \\
\hline
\texttt{CalForce} & \texttt{i}: $\mathbb{R}$, \texttt{ps}: $\mathbb{R}^3$[], 
\texttt{vs}: $\mathbb{R}^3$[], \texttt{info}: FlockingParameter, 
\texttt{goal}: $\mathbb{R}^3$ & $\mathbb{R}^3$ & - \\
\texttt{CalAlignForce} & \texttt{i}: $\mathbb{R}$, \texttt{vs}: $\mathbb{R}^3$[] & $\mathbb{R}^3$ & - \\
\texttt{CalCoheForce} & \texttt{i}: $\mathbb{R}$, \texttt{ps}: $\mathbb{R}^3$[] & $\mathbb{R}^3$ & - \\
\texttt{CalSepaForce} & \texttt{i}: $\mathbb{R}$, \texttt{ps}: $\mathbb{R}^3$[] & $\mathbb{R}^3$ & - \\
\texttt{CalObsAForce} & \texttt{p}: $\mathbb{R}^3$, \texttt{obs}: FlockingParameter.Obs & $\mathbb{R}^3$ & - \\
\texttt{CalGoalForce} & \texttt{p}: $\mathbb{R}^3$, \texttt{v}: $\mathbb{R}^3$, \texttt{goal}: $\mathbb{R}^3$ & $\mathbb{R}^3$ & - \\
\hline
\end{tabular}
\end{center}

\subsection{Semantics}

\subsubsection{Access Routine Semantics}

\noindent CalForce(i, ps, vs, info, goal):
\begin{itemize}
\item output: out := $info.ka \cdot \mathbf{F}_{align} + info.kc \cdot \mathbf{F}_{cohe} + 
info.ks \cdot \mathbf{F}_{sep} + info.ko \cdot \mathbf{F}_{obs} + 
info.kg \cdot \mathbf{F}_{goal} $. 
\end{itemize}

\noindent CalAlignForce(i, vs):
\begin{itemize}
\item output: out := $\mathbf{F}_{align}=\frac{1}{|N_i|}\sum_{j\in N_i}\mathbf{vs}_j-\mathbf{vs}_i$. 
\end{itemize}

\noindent CalCoheForce(i, ps):
\begin{itemize}
\item output: out := $\mathbf{F}_{cohe}=\frac{1}{|N_i|}\sum_{j\in N_i}\mathbf{ps}_j-\mathbf{ps}_i$. 
\end{itemize}

\noindent CalSepaForce(i, ps):
\begin{itemize}
\item output: out := $\mathbf{F}_{sep}=\sum_{j\in N_i}\frac{\mathbf{ps}_i-\mathbf{ps}_j}{\|\mathbf{ps}_i-\mathbf{ps}_j\|^2}$. 
\end{itemize}

\noindent CalObsAForce(p, obs):
\begin{itemize}
\item output: out := $\mathbf{F}_{obs}=\sum_{k\in obs}\frac{\mathbf{p}_i-\mathbf{obs}_k}{\|\mathbf{p}_i-\mathbf{obs}_k\|^2}$. 
\end{itemize}

\noindent CalGoalForce(p, v, goal):
\begin{itemize}
\item output: out := $\mathbf{F}_{goal}=v_{max}\frac{\mathbf{g}-\mathbf{p}_i}{\|\mathbf{g}-\mathbf{p}_i\|}-\mathbf{v}_i$. 
\end{itemize}

~\newpage

\section{MIS of ODE Solver Module} \label{mODE}

\subsection{Module}

ODESolver

\subsection{Uses}

Force

\subsection{Syntax}

\subsubsection{Exported Constants}

None.

\subsubsection{Exported Access Programs}

\begin{center}
\begin{tabular}{p{4cm} p{4cm} p{4cm} p{2cm}}
\hline
\textbf{Name} & \textbf{In} & \textbf{Out} & \textbf{Exceptions} \\
\hline
\texttt{EEulerMethod} & \texttt{ps}: $\mathbb{R}^3$[], 
\texttt{vs}: $\mathbb{R}^3$[], \texttt{info}: FlockingParameter, 
\texttt{goal}: $\mathbb{R}^3$ & $\mathbb{R}^3$[][] & - \\
\texttt{RK4Method} & \texttt{ps}: $\mathbb{R}^3$[], 
\texttt{vs}: $\mathbb{R}^3$[], \texttt{info}: FlockingParameter, 
\texttt{goal}: $\mathbb{R}^3$ & $\mathbb{R}^3$[][] & - \\
\texttt{SIEulerMethod} & \texttt{ps}: $\mathbb{R}^3$[], 
\texttt{vs}: $\mathbb{R}^3$[], \texttt{info}: FlockingParameter, 
\texttt{goal}: $\mathbb{R}^3$& $\mathbb{R}^3$[][] & - \\
\hline
\end{tabular}
\end{center}

\subsection{Semantics}

\subsubsection{Access Routine Semantics}

\noindent EEulerMethod(ps, vs, info, goal):
\begin{itemize}
\item output: \\
a := CalForce(ps, vs, info, goal)
out := $\begin{bmatrix} v_{next} \\ p_{next} \end{bmatrix} = \begin{bmatrix} vs[i] + a \cdot t \\ ps[i] + vs[i] \cdot t \end{bmatrix}$. 
\end{itemize}

\noindent RK4Method(ps, vs, info, goal):
\textit{\# updating the p and v of all agents in this flocking, and solving a = f(i, ps, vs, info, goal) by Force Module for next k}
\begin{itemize}
\item output: \\
$k_1 := \begin{bmatrix} a_1 \\ v \end{bmatrix}$ \\
$k_2 := \begin{bmatrix} a_2 \\ v + \frac{t}{2}k_{1} \end{bmatrix}$ \\
$k_3 := \begin{bmatrix} a_3 \\ v + \frac{t}{2}k_{2} \end{bmatrix}$ \\
$k_4 = \begin{bmatrix} a_4 \\ v + t\,k_{3} \end{bmatrix}$ \\
$\text{out} := \begin{bmatrix} v_{next} \\ p_{next} \end{bmatrix} = \begin{bmatrix} v \\ p \end{bmatrix} + \frac{t}{6}(k_1 + 2k_2 + 2k_3 + k_4)$
\end{itemize}

\noindent SIEulerMethod(ps, vs, info, goal):
\begin{itemize}
\item output: \\
a := CalForce(i, ps, vs, info, goal)
out := $\begin{bmatrix} v_{next} \\ p_{next} \end{bmatrix} = \begin{bmatrix} vs[i] + a \cdot t \\ ps[i] + (vs[i] + a \cdot t) \cdot t \end{bmatrix}$
\end{itemize}

~\newpage

\section{MIS of Flocking Metrics Module} \label{mFM}

\subsection{Module}

FlockingMetrics

\subsection{Uses}

N/A

\subsection{Syntax}

\subsubsection{Exported Access Programs}

\begin{center}
\begin{tabular}{p{4cm} p{4cm} p{4cm} p{2cm}}
\hline
\textbf{Name} & \textbf{In} & \textbf{Out} & \textbf{Exceptions} \\
\hline
\texttt{CalMeanSpeed} & \texttt{vs}: $\mathbb{R}^3$[] & $\mathbb{R}$ & - \\
\texttt{CalAlignMetrics} & \texttt{vs}: $\mathbb{R}^3$[] & $\mathbb{R}$ & - \\
\texttt{CalCenter} & \texttt{ps}: $\mathbb{R}^3$[] & $\mathbb{R}^3$ & - \\
\texttt{CalCohesion} & \texttt{ps}: $\mathbb{R}^3$[], \texttt{center}: $\mathbb{R}^3$ & $\mathbb{R}$ & - \\
\hline
\end{tabular}
\end{center}

\subsection{Semantics}

\subsubsection{Access Routine Semantics}

\noindent CalMeanSpeed(vs):
\begin{itemize}
\item output: out := $\text{avg} = \frac{1}{|vs|} \sum_{i=1}^{|vs|} \left\lVert vs[i] \right\rVert$. 
\end{itemize}

\noindent CalAlignMetrics(vs):
\begin{itemize}
\item output: out := $\text{align} = \frac{1}{|vs|}\left\|\sum_{i=1}^{|vs|}\frac{vs[i]}{\|vs[i]\|}\right\|$. 
\end{itemize}

\noindent CalCenter(ps):
\begin{itemize}
\item output: out := $\text{center} = \frac{1}{|ps|} \sum_{i=1}^{|ps|} ps[i]$. 
\end{itemize}

\noindent CalCohesion(ps, center):
\begin{itemize}
\item output: out := $\text{cohe} =\frac{1}{|ps|} \sum_{i=1}^{|ps|} \left\| ps[i] - \text{center} \right\|$. 
\end{itemize}

~\newpage

\section{MIS of Dynamics Module} \label{mDyna}

\subsection{Module}

Dynamics

\subsection{Uses}

ODESolver, FlockingMetrics

\subsection{Syntax}

\subsubsection{Exported Constants}

\begin{itemize}
\item $goal$ := (500, 500, 0).
\item $arrive$ := 1.
\item $maxTime$ := 1000.
\end{itemize}

\subsubsection{Exported Access Programs}

\begin{center}
\begin{tabular}{p{4cm} p{4cm} p{4cm} p{2cm}}
\hline
\textbf{Name} & \textbf{In} & \textbf{Out} & \textbf{Exceptions} \\
\hline
\texttt{FlockSimulation} & \texttt{initVal}: FlockingParameter, \texttt{method}: str & [AnimationStatus, FlockingStatus] & ErrorMethod \\
\hline
\end{tabular}
\end{center}

\subsection{Semantics}

\subsubsection{State Variables}

\begin{itemize}
\item \texttt{anima}: AnimationStatus, store the position of all agents at every time step. 
\item \texttt{flock}: FlockingStatus, store the flocking metrics at every time step. 
\item \texttt{time}: $\mathbb{R}$, the running time of the flocking simulation. 
\end{itemize}

\subsubsection{Access Routine Semantics}

\noindent FlockSimulation():
\begin{itemize}
\item output: Compute the detail of animation and the metrics of flocking at 
every time steps using ODE Solver Module and Flocking Metrics Module by 
following steps. 
\item exception: exc := ErrorMethod, if the ODE method is not defined.
\end{itemize}
\textit{\# Get the FlockingParameter object from MainController}

\textit{\# Loop until it is timeout}

while not timeout do

\ \ \textit{\# Compute next velocity and position using the selected ODE method}

\ \ [vs\_{next}, ps\_{next}] := ODEChoose()

\ \ \textit{\# Compute next flocking metrics using the FlockingMetrics}

\ \ speed := CalMeanSpeed()

\ \ align := CalAlignMetrics()

\ \ cohe := CalCohesion()

\ \ center := CalCenter()

\ \ \textit{\# Store the agent position in anima}

\ \ anima.add(ps\_{next})

\ \ \textit{\# Store the flocking metrics in flock}

\ \ flock.add(speed, align, cohe)

\ \ \textit{\# Check if the flocking reach the goal}

\ \ if (distance between goal and center) $\leq$ arrive, then break

\ \ \textit{\# Update the simulation time}

end while

\textit{\# Return [anima, flock]}

\subsubsection{Local Functions}

\noindent ODEChoose():
\begin{itemize}
\item input: \texttt{method}: str, \texttt{i}: $\mathbb{R}$, \texttt{ps}: $\mathbb{R}^3$[], 
\texttt{vs}: $\mathbb{R}^3$[], \texttt{info}: FlockingParameter. 
\item output: $\mathbb{R}^3$[].
\item exception: exc := ErrorMethod, if the ODE method is not defined.
\item description: Compute the velocity and position at next time step by 
specific method in ODESolver Module.
\end{itemize}

~\newpage

\section{MIS of Flocking Parameter Module} \label{mFP}

\subsection{Module}

FlockingParameter

\subsection{Uses}

N/A

\subsection{Syntax}

The class uses the default constructor provided by the language. All 
attributes are publicly accessible.

\subsubsection{Exported Types}

\begin{itemize}
\item \texttt{Obs}: the structure of obstacle.
\begin{itemize}
  \item \texttt{p}: $\mathbb{R}^3$, the center of circle obstacle.
  \item \texttt{r}: $\mathbb{R}$, the radius of circle obstacle.
\end{itemize}
\end{itemize}

\subsection{Semantics}

\subsubsection{State Variables}

\begin{itemize}
\item \texttt{ka}: $\mathbb{R}$, the weight of alignment force. 
\item \texttt{kc}: $\mathbb{R}$, the weight of cohesion force. 
\item \texttt{ks}: $\mathbb{R}$, the weight of separation force. 
\item \texttt{ko}: $\mathbb{R}$, the weight of obstacle avoidance force. 
\item \texttt{kg}: $\mathbb{R}$, the weight of goal-seeking force. 
\item \texttt{n}: $\mathbb{Z}$, the number of agent. 
\item \texttt{obs}: array of Obs, the information of obstacles;
\item \texttt{t}: $\mathbb{R}$, the length of time step. 
\item \texttt{starts}: array of $\mathbb{R}^3$, the start points of agents. 
\end{itemize}

~\newpage

\section{MIS of Animation Status Module} \label{mAS}

\subsection{Module}

AnimationStatus

\subsection{Uses}

N/A

\subsection{Syntax}

The class uses the default constructor provided by the language. All 
attributes are publicly accessible.

\subsection{Semantics}

\subsubsection{State Variables}

\begin{itemize}
\item \texttt{t}: $\mathbb{R}$, the length of time step. 
\item \texttt{agents}: ($\mathbb{R}^3$[])[], the position of agents. 
\item \texttt{obs}: array of Obs, the information of obstacles;
\end{itemize}


~\newpage

\section{MIS of Flocking Status Module} \label{mFS}

\subsection{Module}

FlockingStatus

\subsection{Uses}

N/A

\subsection{Syntax}

The class uses the default constructor provided by the language. All 
attributes are publicly accessible.

\subsection{Semantics}

\subsubsection{State Variables}

\begin{itemize}
\item \texttt{t}: $\mathbb{R}$, the length of time step. 
\item \texttt{meanSpeed}: $\mathbb{R}$[], the list of mean speed at every time step. 
\item \texttt{alignMetric}: $\mathbb{R}$[], the list of align at every time step.
\item \texttt{cohesion}: $\mathbb{R}$[], the list of cohesion at every time step.
\end{itemize}

~\newpage

\section{MIS of Input Reader Module} \label{mInput}

\subsection{Module}

InputReader

\subsection{Uses}

Hardware-Hiding, FilesIO, InputValidation

\subsection{Syntax}

\subsubsection{Exported Constants}

None.

\subsubsection{Exported Access Programs}

\begin{center}
\begin{tabular}{p{4cm} p{4cm} p{4cm} p{2cm}}
\hline
\textbf{Name} & \textbf{In} & \textbf{Out} & \textbf{Exceptions} \\
\hline
\texttt{GetUserInput} & - & FlockingParameter & MissingK, 
MissingN, MissingT, MissingP, MissingR, ErrorK, ErrorN, ErrorT, ErrorP, ErrorR \\ \\
\texttt{GetSimulateParam} & - & FlockingParameter & ErrorObject \\
\texttt{GetMethod} & - & str & MissingMethod, ErrorMethod \\
\texttt{GetAnimaStatus} & - & AnimationStatus & ErrorObject \\
\texttt{GetFlockStatus} & - & FlockingStatus[] & ErrorObject \\
\hline
\end{tabular}
\end{center}

\subsection{Semantics}

\subsubsection{Environment Variables}

\begin{itemize}
\item kaInput: the UI input object used to read the value of parameter ka.
\item kcInput: the UI input object used to read the value of parameter kc.
\item ksInput: the UI input object used to read the value of parameter ks.
\item koInput: the UI input object used to read the value of parameter ko.
\item kgInput: the UI input object used to read the value of parameter kg.
\item tInput: the UI input object used to read the value of length of timestep.
\item nInput: the UI input object used to read the value of number of agent.
\item obsInput: the UI input object used to read the information of obstacles.
\item methodInput: the UI input object used to read the user input ODE method.
\item initPathInput: the UI input object used to read the file path of initial flocking parameters.
\item animaPathInput: the UI input object used to read the file path of animation status.
\item flockPathInput: the UI input object used to read the file path of flocking metrics.
\end{itemize}

\subsubsection{Access Routine Semantics}

\noindent GetUserInput():
\begin{itemize}
\item output: \texttt{out} := FlockingParameter, a new FlockingParameter object 
defined by the user-provided values
\item exception: all exceptions are thrown by InputValidation.
% \item output: out := (get the parameters user input and set them into a new FlockingParameter)
% \item exception: exc := ((missing any parameters in out.ka, out.kc, out.ks, 
% out.ko, out.kg, out.n, out.t )=> MissingInput | (missing p or r at any out.obs[i]) => ErrorObs)
\end{itemize}

\noindent GetSimulateParam():
\begin{itemize}
\item output: \texttt{out} := FlockingParameter, a new FlockingParameter object 
read from the user-specified path.
\item exception: 
\begin{itemize}
    \item ErrorObject, if there is no such file or the FlockingParameter read 
    from the file is null
\end{itemize}
\end{itemize}

\noindent GetMethod():
\begin{itemize}
\item output: \texttt{out} := str, the ODE method user chooses.
\item exception: all exceptions are thrown by InputValidation.
\end{itemize}

\noindent GetAnimaStatus():
\begin{itemize}
\item output: \texttt{out} := AnimationStatus[], a list of AnimationStatus object 
read from the user-specified paths.
\item exception: 
\begin{itemize}
    \item ErrorObject, if for any path, there is no such file or the 
    AnimationStatus read from the file is null
\end{itemize}
\end{itemize}

\noindent GetFlockStatus():
\begin{itemize}
\item output: \texttt{out} := FlockingStatus[], a list of FlockingStatus object 
read from the user-specified paths.
\item exception: 
\begin{itemize}
    \item ErrorObject, if for any path, there is no such file or the 
    FlockingStatus read from the file is null
\end{itemize}
\end{itemize}

~\newpage

\section{MIS of File I/O Module} \label{mFile}

\subsection{Module}

FilesIO

\subsection{Uses}

N/A

\subsection{Syntax}

\subsubsection{Exported Constants}

None.

\subsubsection{Exported Access Programs}

\begin{center}
\begin{tabular}{p{4cm} p{4cm} p{4cm} p{2cm}}
\hline
\textbf{Name} & \textbf{In} & \textbf{Out} & \textbf{Exceptions} \\
\hline
\texttt{Read} & \texttt{path}: str, \texttt{T}: type & T & NoFileFound \\
\texttt{Write} & \texttt{path}: str, \texttt{contents}: str & - & - \\
\hline
\end{tabular}
\end{center}

\subsection{Semantics}

\subsubsection{Access Routine Semantics}

\noindent Read(path):
\begin{itemize}
\item output: out := the contents of file.
\item exception: 
\begin{itemize}
    \item NoFileFound, if the file at specific path is not found.
\end{itemize}
\end{itemize}

\noindent Write(path, contents):
\begin{itemize}
\item output: None.
\end{itemize}

~\newpage

\section{MIS of Animation Simulator Module} \label{mAnima}

\subsection{Module}

AnimationSimulator

\subsection{Uses}

N/A

\subsection{Syntax}

\subsubsection{Exported Constants}

None.

\subsubsection{Exported Access Programs}

\begin{center}
\begin{tabular}{p{4cm} p{4cm} p{4cm} p{2cm}}
\hline
\textbf{Name} & \textbf{In} & \textbf{Out} & \textbf{Exceptions} \\
\hline
\texttt{PlaySimulation} & \texttt{anima}: AnimationStatus & - & - \\
\hline
\end{tabular}
\end{center}

\subsection{Semantics}

\subsubsection{Access Routine Semantics}

\noindent PlaySimulation(anima):
\begin{itemize}
\item output: None, the animation will be displayed in the Application Window.
\end{itemize}

~\newpage

\section{MIS of Chart Drawer Module} \label{mChart}

\subsection{Module}

ChartDrawer

\subsection{Uses}

N/A

\subsection{Syntax}

\subsubsection{Exported Constants}

None.

\subsubsection{Exported Access Programs}

\begin{center}
\begin{tabular}{p{4cm} p{4cm} p{4cm} p{2cm}}
\hline
\textbf{Name} & \textbf{In} & \textbf{Out} & \textbf{Exceptions} \\
\hline
\texttt{DrawChart} & \texttt{metrics}: FlockingStatus[] & - & - \\
\hline
\end{tabular}
\end{center}

\subsection{Semantics}

\subsubsection{Access Routine Semantics}

\noindent DrawChart(metrics):
\begin{itemize}
\item output: None, and the line charts will be displayed in the Application Window.
\end{itemize}

% ~\newpage



% \section{MIS of \wss{Module Name}} \label{Module} \wss{Use labels for
%   cross-referencing}

% \wss{You can reference SRS labels, such as R\ref{R_Inputs}.}

% \wss{It is also possible to use \LaTeX for hypperlinks to external documents.}

% \subsection{Module}

% \wss{Short name for the module}

% \subsection{Uses}


% \subsection{Syntax}

% \subsubsection{Exported Constants}

% \subsubsection{Exported Access Programs}

% \begin{center}
% \begin{tabular}{p{2cm} p{4cm} p{4cm} p{2cm}}
% \hline
% \textbf{Name} & \textbf{In} & \textbf{Out} & \textbf{Exceptions} \\
% \hline
% \wss{accessProg} & - & - & - \\
% \hline
% \end{tabular}
% \end{center}

% \subsection{Semantics}

% \subsubsection{State Variables}

% \wss{Not all modules will have state variables.  State variables give the module
%   a memory.}

% \subsubsection{Environment Variables}

% \wss{This section is not necessary for all modules.  Its purpose is to capture
%   when the module has external interaction with the environment, such as for a
%   device driver, screen interface, keyboard, file, etc.}

% \subsubsection{Assumptions}

% \wss{Try to minimize assumptions and anticipate programmer errors via
%   exceptions, but for practical purposes assumptions are sometimes appropriate.}

% \subsubsection{Access Routine Semantics}

% \noindent \wss{accessProg}():
% \begin{itemize}
% \item transition: \wss{if appropriate} 
% \item output: \wss{if appropriate} 
% \item exception: \wss{if appropriate} 
% \end{itemize}

% \wss{A module without environment variables or state variables is unlikely to
%   have a state transition.  In this case a state transition can only occur if
%   the module is changing the state of another module.}

% \wss{Modules rarely have both a transition and an output.  In most cases you
%   will have one or the other.}

% \subsubsection{Local Functions}

% \wss{As appropriate} \wss{These functions are for the purpose of specification.
%   They are not necessarily something that is going to be implemented
%   explicitly.  Even if they are implemented, they are not exported; they only
%   have local scope.}

\newpage

\bibliographystyle {plainnat}
\bibliography {../../../refs/References}

% \newpage

% \section{Appendix} \label{Appendix}

% \wss{Extra information if required}

% \newpage{}

% \section*{Appendix --- Reflection}

% \wss{Not required for CAS 741 projects}

% The information in this section will be used to evaluate the team members on the
% graduate attribute of Problem Analysis and Design.

% \input{../../Reflection.text}

% \begin{enumerate}
%   \item What went well while writing this deliverable? 
%   \item What pain points did you experience during this deliverable, and how
%     did you resolve them?
%   \item Which of your design decisions stemmed from speaking to your client(s)
%   or a proxy (e.g. your peers, stakeholders, potential users)? For those that
%   were not, why, and where did they come from?
%   \item While creating the design doc, what parts of your other documents (e.g.
%   requirements, hazard analysis, etc), it any, needed to be changed, and why?
%   \item What are the limitations of your solution?  Put another way, given
%   unlimited resources, what could you do to make the project better? (LO\_ProbSolutions)
%   \item Give a brief overview of other design solutions you considered.  What
%   are the benefits and tradeoffs of those other designs compared with the chosen
%   design?  From all the potential options, why did you select the documented design?
%   (LO\_Explores)
%   \item (After you have implemented another team's module, which means this
%   isn't filled in until after the original deadline). What did you learn by
%   implementing another team's module? Were all the details you needed in the
%   documentation, or did you need to make assumptions, or ask the other team
%   questions?  If your team also had another team implement one of your modules,
%   what was this experience like?  Are there things in your documentation you
%   could have changed to make the process go more smoothly for when an
%   ``outsider'' completes some of the implementation?
% \end{enumerate}


\end{document}